% !TeX root = ../algebra-1.tex
\chapter{Subspaces and dimension}\label{chap:algebra-subspaces}

\section{Subspaces, kernels and column spaces}

\begin{definition}
A subset $U\subseteq F^n$ is a \emph{subspace} if it is nonempty and is
closed under addition and scalar multiplication:
\[
 u,v\in U,\ a\in F \quad\Longrightarrow\quad u+v\in U,\ au\in U.
\]
We write $U\leq F^n$. More generally, $U\leq V$ means that $U$ is a
subspace contained in the subspace $V$. It is \emph{proper} if $U\ne V$.
\end{definition}

Equivalently, a nonempty subset is a subspace if $au+bv$ belongs to it
whenever $u,v$ do and $a,b\in F$.
The scalar field matters: closure must hold for every scalar in the stated
field, including negative scalars and, over $\mathbb C$, complex scalars.

\begin{proposition}
Every subspace contains the zero vector.
\end{proposition}
\begin{proof}
Choose $u\in U$, using nonemptiness. Closure under multiplication by $0$
gives $0u=0_n\in U$.
\end{proof}

\begin{example}
The sets $\{0_n\}$ and $F^n$ are subspaces. In $\mathbb R^2$, a line
through the origin, such as $\{(t,2t)^{\mathsf T}:t\in\mathbb R\}$, is a
subspace. The half-line with $t\geq0$ is not: multiplication by $-1$
takes its nonzero vectors outside the set.
A line not through the origin fails the zero-vector condition.
\end{example}

\begin{remark}
Containing zero is necessary, but does not by itself prove that a set is a
subspace. For instance, the union of the two coordinate axes contains
zero but does not contain $(1,0)^{\mathsf T}+(0,1)^{\mathsf T}$.
\end{remark}

\begin{definition}
For an $m\times n$ matrix $A$, its \emph{kernel} and \emph{column space} are
\[
 N(A)=\{v\in F^n:Av=0_m\},\qquad
 C(A)=\{Av:v\in F^n\}\subseteq F^m.
\]
The kernel describes inputs mapped to zero; the column space describes
all possible outputs.
\end{definition}

\begin{lemma}[Kernels and column spaces are subspaces]
$N(A)\leq F^n$ and $C(A)\leq F^m$.
\end{lemma}
\begin{proof}
Since $A0_n=0_m$, both sets are nonempty.
If $u,v\in N(A)$ and $a,b\in F$, then
\[
 A(au+bv)=aAu+bAv=a0_m+b0_m=0_m,
\]
so $au+bv\in N(A)$. If $y=Au$ and $z=Av$ belong to $C(A)$, then
\[
 ay+bz=aAu+bAv=A(au+bv)\in C(A).
\]
Thus each set has the required closure properties.
\end{proof}

\section{Spanning sequences and linear independence}

\begin{definition}
The \emph{span} of vectors $v_1,\ldots,v_k$ is
\[
 \spanop(v_1,\ldots,v_k)
 =\{a_1v_1+\cdots+a_kv_k:a_1,\ldots,a_k\in F\}.
\]
They \emph{span} $U$ if this set equals $U$. The span of the empty
sequence is defined to be $\{0_n\}$.
\end{definition}

\begin{proposition}
The span is a subspace and is the smallest subspace containing all the
listed vectors. The column space of $A=[\,c_1\ \cdots\ c_n\,]$ is
$\spanop(c_1,\ldots,c_n)$.
\end{proposition}
\begin{proof}
Zero is a linear combination with all coefficients zero. Sums and scalar
multiples of linear combinations are again linear combinations.
Any subspace containing the listed vectors contains all their linear
combinations, which proves minimality. Finally,
$Ax=x_1c_1+\cdots+x_nc_n$, so the possible outputs are exactly their span.
\end{proof}

\begin{example}
The $xy$-plane in $\mathbb R^3$ is spanned by $e_1,e_2$, and also by
$(1,1,0)^{\mathsf T},(1,-1,0)^{\mathsf T}$, since
\[
 (a,b,0)^{\mathsf T}
 =\frac{a+b}{2}(1,1,0)^{\mathsf T}
  +\frac{a-b}{2}(1,-1,0)^{\mathsf T}.
\]
The sequence $e_1,e_2,e_3$ spans all of $\mathbb R^3$, so it does not
span this plane: its span is too large.
\end{example}

\begin{example}[Testing a span by solving a system]
Let
\[
 u=\begin{pmatrix}1\\2\\3\end{pmatrix},\quad
 v=\begin{pmatrix}4\\5\\6\end{pmatrix},\quad
 w=\begin{pmatrix}7\\8\\9\end{pmatrix}.
\]
To decide whether $(a,b,c)^{\mathsf T}$ belongs to their span, solve the
system with coefficient matrix $[\,u\ v\ w\,]$. Row reduction gives
\[
 \left[\begin{array}{ccc|c}
 1&4&7&a\\2&5&8&b\\3&6&9&c
 \end{array}\right]
 \ \longrightarrow\
 \left[\begin{array}{ccc|c}
 1&0&-1&(4b-5a)/3\\
 0&1&2&(2a-b)/3\\
 0&0&0&a-2b+c
 \end{array}\right].
\]
Consequently,
\[
 \spanop(u,v,w)
 =\{(a,b,c)^{\mathsf T}:a-2b+c=0\}.
\]
This is a plane, not all of $\mathbb R^3$.
\end{example}

\begin{definition}
A sequence $v_1,\ldots,v_k$ is \emph{linearly independent} if
\[
 a_1v_1+\cdots+a_kv_k=0_n
 \quad\Longrightarrow\quad a_1=\cdots=a_k=0.
\]
Otherwise it is \emph{linearly dependent}. The empty sequence is independent.
\end{definition}

A sequence containing zero is dependent. A sequence with a repeated vector
is dependent. More generally, dependence means that at least one listed
vector can be expressed as a linear combination of the others.

\begin{example}
The vectors $(1,0)^{\mathsf T},(0,1)^{\mathsf T},(1,1)^{\mathsf T}$
are dependent because the first two minus the third give zero.
By contrast, $(1,1)^{\mathsf T},(1,-1)^{\mathsf T}$ are independent:
\[
 a(1,1)^{\mathsf T}+b(1,-1)^{\mathsf T}=0
 \quad\Longrightarrow\quad a+b=a-b=0
 \quad\Longrightarrow\quad a=b=0.
\]
For the preceding three-dimensional example, $u-2v+w=0$, so $u,v,w$
are dependent.
\end{example}

\method{Put the vectors in the columns of a matrix $A$. They are
independent exactly when $Ax=0$ has only the zero solution.
To test whether they span $F^m$, ask whether $Ax=b$ is solvable for
every $b\in F^m$. These are different questions.}

\section{Bases and dimension}

\begin{definition}
A \emph{basis} of $U$ is an ordered sequence of vectors that is both
linearly independent and spans $U$.
\end{definition}

Spanning ensures that every vector has a representation in the basis.
Independence ensures that this representation is unique: subtract two
representations and apply the definition of independence.
The standard basis of $F^n$ is $e_1,\ldots,e_n$.

\begin{lemma}[Steinitz length bound; source Lemma 2.15]
Suppose $s_1,\ldots,s_m$ span $U$. Any sequence of more than $m$ vectors
in $U$ is linearly dependent.
\end{lemma}
\begin{proof}
Let $v_1,\ldots,v_k\in U$ with $k>m$. Write
$v_j=\sum_{i=1}^m a_{ij}s_i$, and form the $m\times k$ matrix
$A=(a_{ij})$. The homogeneous system $Ax=0$ has more unknowns than
equations, so it has a nonzero solution $x$. Then
\[
 \sum_{j=1}^k x_jv_j
 =\sum_{i=1}^m\left(\sum_{j=1}^k a_{ij}x_j\right)s_i=0.
\]
Since some $x_j$ is nonzero, this is a nontrivial dependence.
The case $m=0$ simply says that every vector of $U=\{0\}$ is zero.
\end{proof}

\begin{corollary}
Any $n+1$ vectors in $F^n$ are dependent.
\end{corollary}
\begin{proof}
Apply the lemma to the $n$ standard basis vectors.
\end{proof}

\begin{theorem}[Size of a basis; source Theorem 2.17]
Any two bases of a subspace have the same number of vectors.
\end{theorem}
\begin{proof}
Let one basis have $r$ vectors and the other $s$.
The first spans the subspace and the second is independent, so the
Steinitz bound gives $s\leq r$. Reversing the roles gives $r\leq s$.
Thus $r=s$.
\end{proof}

\begin{lemma}[Extending an independent sequence]
If $u_1,\ldots,u_k$ are independent and
$v\notin\spanop(u_1,\ldots,u_k)$, then $u_1,\ldots,u_k,v$ are independent.
\end{lemma}
\begin{proof}
If $\sum_i a_i u_i+bv=0$ and $b\ne0$, then
$v=-b^{-1}\sum_i a_i u_i$, contrary to the assumption.
Thus $b=0$, after which independence of the $u_i$ gives all $a_i=0$.
\end{proof}

\begin{theorem}[Basis extension]
Every independent sequence in a subspace $U\leq F^n$ can be extended to
a basis of $U$. In particular, every subspace has a basis.
\end{theorem}
\begin{proof}
If the current sequence does not span $U$, choose a vector of $U$ outside
its span and append it. The lemma preserves independence.
This process cannot produce more than $n$ independent vectors, so it
must stop at a spanning sequence. Starting with the empty sequence gives
existence, including the empty basis of $\{0_n\}$.
\end{proof}

\begin{definition}
The \emph{dimension} of $U$, written $\dim U$, is the number of vectors
in any basis of $U$. The size theorem makes this well-defined.
\end{definition}

\begin{example}
$\dim F^n=n$ and $\dim\{0_n\}=0$.
The plane $a+b+c=0$ has basis
$(1,0,-1)^{\mathsf T},(0,1,-1)^{\mathsf T}$, because
$(a,b,c)^{\mathsf T}=a(1,0,-1)^{\mathsf T}
+b(0,1,-1)^{\mathsf T}$ whenever $c=-a-b$.
Its dimension is two.
\end{example}

\begin{proposition}
If $U\leq V\leq F^n$, then $\dim U\leq\dim V$.
Equality holds if and only if $U=V$.
\end{proposition}
\begin{proof}
Extend a basis of $U$ to a basis of $V$. This can only increase its
length. If the dimensions are equal, no extra vectors are needed, so
the original basis already spans $V$ and $U=V$.
\end{proof}

\begin{corollary}
In a $d$-dimensional subspace, any $d$ independent vectors form a basis.
Also, any $d$ vectors that span it form a basis.
\end{corollary}
\begin{proof}
For the first assertion, basis extension leaves no room for extra vectors.
For the second, if a spanning sequence is dependent, delete a vector
expressible using the others. Repeating gives a basis of length less
than $d$, a contradiction.
\end{proof}

\section{Sums and intersections}

\begin{definition}
For subspaces $S,T\leq F^n$, define
\[
 S\cap T=\{v:v\in S\text{ and }v\in T\},\qquad
 S+T=\{s+t:s\in S,\ t\in T\}.
\]
\end{definition}

Both are subspaces. For the intersection, a linear combination stays in
each of the two sets. For the sum,
$a(s_1+t_1)+b(s_2+t_2)=(as_1+bs_2)+(at_1+bt_2)\in S+T$.
Both sets contain zero.
The sum is the span of the combined spanning sequences of $S$ and $T$.
It usually differs from the union $S\cup T$.

\begin{example}
Let $S=\spanop((2,-1,0)^{\mathsf T})$ and
$T=\spanop((2,-1,2)^{\mathsf T})$ in $\mathbb R^3$.
The two generating vectors are independent, so their sum is a plane.
If $a(2,-1,0)^{\mathsf T}=b(2,-1,2)^{\mathsf T}$, the third coordinate
gives $b=0$, and then $a=0$. Thus $S\cap T=\{0_3\}$.
\end{example}

\begin{theorem}[Dimension formula; source Theorem 2.28]
\[
 \boxed{\dim(S+T)=\dim S+\dim T-\dim(S\cap T).}
\]
\end{theorem}
\begin{proof}
Choose a basis $v_1,\ldots,v_k$ of $S\cap T$.
Extend it to a basis $v_1,\ldots,v_k,s_1,\ldots,s_r$ of $S$, and
to a basis $v_1,\ldots,v_k,t_1,\ldots,t_q$ of $T$.
The combined sequence $v_1,\ldots,v_k,s_1,\ldots,s_r,t_1,\ldots,t_q$
spans $S+T$.

For independence, suppose
$\sum_i a_iv_i+\sum_j b_js_j+\sum_\ell c_\ell t_\ell=0$.
Then $\sum_\ell c_\ell t_\ell\in S\cap T$, so it is some
$\sum_i d_iv_i$. Independence of the chosen basis of $T$ forces
all $c_\ell$ and $d_i$ to be zero.
Independence of the basis of $S$ then gives all $a_i=b_j=0$.
Thus the combined sequence is a basis of length $k+r+q$, which equals
$(k+r)+(k+q)-k$.
\end{proof}

For example, two distinct planes through the origin in $\mathbb R^3$
intersect in a line and have sum $\mathbb R^3$.
If $S\cap T=\{0\}$, the sum is called a \emph{direct sum}, written
$S\oplus T$; the representation $s+t$ is then unique.

\section{Computing bases from a matrix}

\subsection{A basis for the kernel}

For an RREF matrix, assign each free variable a parameter.
Set one free variable to $1$ and the others to $0$, then solve for the
pivot variables. The resulting vectors are the \emph{fundamental solutions}.

\begin{example}[Source Example 2.30]
Suppose $A$ has reduced row echelon form
\[
 R=\begin{pmatrix}0&1&2&0&3\\0&0&0&1&4\\0&0&0&0&0\end{pmatrix}.
\]
The equations are $x_2=-2x_3-3x_5$ and $x_4=-4x_5$.
Variables $x_1,x_3,x_5$ are free, so
\[
 N(A)=N(R)=\spanop\left(
 \begin{pmatrix}1\\0\\0\\0\\0\end{pmatrix},
 \begin{pmatrix}0\\-2\\1\\0\\0\end{pmatrix},
 \begin{pmatrix}0\\-3\\0\\-4\\1\end{pmatrix}\right).
\]
These vectors are independent: their first, third and fifth coordinates
recover their coefficients. They therefore form a basis, and $\dim N(A)=3$.
\end{example}

\begin{proposition}
The fundamental solutions form a basis of $N(A)$. Hence $\dim N(A)$
is the number of free variables.
\end{proposition}
\begin{proof}
Every solution is the linear combination whose coefficients are its
free-variable values, so the fundamental solutions span.
At the free-variable positions they have the columns of an identity
matrix, so a zero linear combination must have every coefficient zero.
\end{proof}

\subsection{A basis for the column space}

\begin{definition}
The \emph{column rank}, or simply \emph{rank}, is
$\rank A=\dim C(A)$.
\end{definition}

\begin{theorem}
If the pivot columns of the RREF of $A$ have indices
$j_1,\ldots,j_r$, then columns $j_1,\ldots,j_r$ of the
\textbf{original matrix $A$} form a basis of $C(A)$.
\end{theorem}
\begin{proof}
Write $R=EA$, where $E$ is the invertible product of the elementary
matrices used in row reduction. For every coefficient vector $x$,
$Ax=0$ if and only if $Rx=0$. Thus the columns of $A$ and $R$ have
exactly the same linear relations.

In $R$, the pivot columns are the first $r$ standard coordinate vectors.
Every other column is their linear combination.
The identical relations hold in $A$, so its corresponding columns span
$C(A)$. Any dependence among these columns would give a dependence among
the pivot columns of $R$, which are independent.
\end{proof}

\begin{remark}
Row operations preserve the kernel and the relations between columns,
but can change the column space itself. For example,
\[
 A=\begin{pmatrix}1&0\\1&0\end{pmatrix}
 \longrightarrow R=\begin{pmatrix}1&0\\0&0\end{pmatrix}.
\]
Here $C(A)=\spanop((1,1)^{\mathsf T})$ and
$C(R)=\spanop((1,0)^{\mathsf T})$ are different.
\end{remark}

\begin{example}
The matrix in Example~\ref{ex:algebra-rref} has pivots in columns one and
two. Therefore a column-space basis is
$(1,4,7)^{\mathsf T},(2,5,8)^{\mathsf T}$, and its rank is two.
\end{example}

\begin{theorem}[Rank--nullity; source Theorem 2.39]
For every $m\times n$ matrix $A$,
\[
 \boxed{\dim N(A)+\rank A=n.}
\]
\end{theorem}
\begin{proof}
There are $n$ columns. Each is either a pivot column or a free-variable
column. The two numbers are $\rank A$ and $\dim N(A)$ respectively.
\end{proof}

\subsection{A basis for the row space}

Representing rows as column vectors, the \emph{row space} of $A$ is
$C(A^{\mathsf T})\leq F^n$.

\begin{lemma}
If $F_0$ is an invertible $n\times n$ matrix, then $C(AF_0)=C(A)$.
\end{lemma}
\begin{proof}
Every $AF_0x$ is in $C(A)$. Conversely, every $Ay$ equals
$AF_0(F_0^{-1}y)$, so it lies in $C(AF_0)$.
\end{proof}

Consequently, row operations preserve row space:
if $R=EA$, then $R^{\mathsf T}=A^{\mathsf T}E^{\mathsf T}$ and the
lemma applies. The nonzero rows of $R$, transposed, form a row-space
basis. They span by definition, and their pivot positions prove
independence.

\begin{theorem}[Equality of row and column rank]
\[
 \dim C(A^{\mathsf T})=\dim C(A)=\rank A.
\]
\end{theorem}
\begin{proof}
Both dimensions equal the number of pivots in an echelon form.
\end{proof}

Equal dimensions do not mean equal spaces. For
$A=\begin{pmatrix}1&1\\2&2\end{pmatrix}$, the column space is
$\spanop((1,2)^{\mathsf T})$, while the row space is
$\spanop((1,1)^{\mathsf T})$.

\begin{center}
\begin{tabular}{lll}
\toprule
Space & Ambient space & Basis obtained from RREF $R$\\
\midrule
$N(A)$ & $F^n$ & Fundamental solutions of $Rx=0$\\
$C(A)$ & $F^m$ & Original columns at pivot indices\\
$C(A^{\mathsf T})$ & $F^n$ & Nonzero rows of $R$, transposed\\
\bottomrule
\end{tabular}
\end{center}

\section{Equivalent tests for invertibility}

\begin{theorem}[Big invertibility theorem; source Theorem 2.40]
For a square $n\times n$ matrix $A$, the following are equivalent:
\begin{enumerate}
\item $A$ is invertible.
\item $A^{\mathsf T}$ is invertible.
\item The RREF of $A$ is $I_n$.
\item $N(A)=\{0_n\}$.
\item The columns of $A$ are independent.
\item $\rank A=n$.
\item The columns of $A$ form a basis of $F^n$.
\item $C(A)=F^n$.
\item For every $b\in F^n$, the equation $Ax=b$ has a unique solution.
\end{enumerate}
\end{theorem}
\begin{proof}
The inverse criterion from Chapter~\ref{chap:algebra-matrices} equates
invertibility with RREF $I_n$.
For a square matrix, this means there are $n$ pivots, equivalently no
free variables, equivalently a trivial kernel.
A trivial kernel is precisely independence of the columns.
The column-space basis theorem gives rank equal to the number of pivots.
In dimension $n$, $n$ independent columns form a basis, equivalently they
span $F^n$. In that basis, every $b$ has a unique coefficient vector
$x$, which is the unique solution of $Ax=b$.
Conversely, unique solvability at $b=0$ gives a trivial kernel.
Finally, if $A$ is invertible, then
$(A^{\mathsf T})^{-1}=(A^{-1})^{\mathsf T}$; applying the same argument
to $A^{\mathsf T}$ gives the reverse implication.
\end{proof}

The determinant criterion $\det A\ne0$ will be added in
Chapter~\ref{chap:algebra-eigenvalues}.
For a rectangular matrix, keep the distinctions:
independent columns mean rank $n$, whereas spanning the codomain $F^m$
means rank $m$. They coincide only when $m=n$.
